\documentclass[11pt]{article}
\usepackage[T1]{fontenc}
\usepackage[utf8]{inputenc}
\usepackage{amsmath,amsthm,mathtools}
\usepackage{newpxtext,newpxmath}
\usepackage[a4paper,margin=26mm,headheight=15pt]{geometry}
\usepackage{microtype}
\usepackage{booktabs,tabularx,array}
\usepackage{xcolor}
\usepackage[most]{tcolorbox}
\usepackage{enumitem}
\usepackage{titlesec,fancyhdr}
\usepackage{listings}
\usepackage{hyperref}
\definecolor{Ink}{HTML}{233C53}
\definecolor{Accent}{HTML}{1F6C6A}
\definecolor{Pale}{HTML}{F0F5F6}
\definecolor{Muted}{HTML}{526370}
\hypersetup{colorlinks=true,linkcolor=Accent,citecolor=Accent,urlcolor=Accent,
  pdftitle={A Reflexive Root-Polytope Model for Preorder h-Polynomials},
  pdfauthor={Research manuscript prepared with ChatGPT},
  pdfsubject={Proposed proof of the simplicial-polytopal realization conjecture for finite preorders}}
\titleformat{\section}{\Large\bfseries\color{Ink}}{\thesection}{0.7em}{}
\titleformat{\subsection}{\large\bfseries\color{Ink}}{\thesubsection}{0.7em}{}
\pagestyle{fancy}
\fancyhf{}
\fancyhead[L]{\small\textsc{Preorder $h$-polynomials}}
\fancyhead[R]{\small\textsc{Research manuscript}}
\fancyfoot[C]{\thepage}
\renewcommand{\headrulewidth}{0.3pt}
\setlength{\parskip}{4pt}
\setlength{\parindent}{1.1em}
\setlength{\emergencystretch}{2em}
\setcounter{tocdepth}{1}
\numberwithin{equation}{section}
\newtheorem{theorem}{Theorem}[section]
\newtheorem{lemma}[theorem]{Lemma}
\newtheorem{proposition}[theorem]{Proposition}
\newtheorem{corollary}[theorem]{Corollary}
\theoremstyle{definition}
\newtheorem{definition}[theorem]{Definition}
\newtheorem{remark}[theorem]{Remark}
\newcommand{\R}{\mathbb R}
\newcommand{\Z}{\mathbb Z}
\newcommand{\N}{\mathbb Z_{\geq0}}
\newcommand{\J}{\mathcal J(\tau)}
\newcommand{\Qt}{\mathcal Q_\tau}
\newcommand{\Pt}{P_\tau}
\newcommand{\Bt}{B_\tau}
\newcommand{\Ct}{C_\tau}
\newcommand{\hpoly}{h_\tau(t)}
\newcommand{\conv}{\operatorname{conv}}
\newcommand{\supp}{\operatorname{supp}}
\newcommand{\Ehr}{\operatorname{Ehr}}
\newcommand{\relint}{\operatorname{relint}}
\newcommand{\Vol}{\operatorname{Vol}_{\Z}}
\newcommand{\one}{\mathbf 1}
\newcommand{\zero}{\mathbf 0}
\newcommand{\ee}{\mathbf e}
\newcommand{\ff}{\mathbf f}
\newcommand{\down}{\mathord\downarrow}
\newcommand{\arxiv}[1]{\href{https://arxiv.org/abs/#1}{arXiv:#1}}
\newcolumntype{Y}{>{\raggedright\arraybackslash}X}
\lstset{basicstyle=\small\ttfamily,breaklines=true,columns=fullflexible,
  backgroundcolor=\color{Pale},frame=single,rulecolor=\color{Pale},
  xleftmargin=0.5em,xrightmargin=0.5em,aboveskip=8pt,belowskip=8pt}

\begin{document}
\thispagestyle{empty}
\begin{flushleft}
{\small\sffamily\bfseries\color{Accent} ORDER THEORY \ / \ LATTICE-POINT ENUMERATION}\par
\vspace{10pt}
{\LARGE\bfseries\color{Ink} A Reflexive Root-Polytope Model\\[4pt]
for Preorder $h$-Polynomials}\par
\vspace{9pt}
{\large Simplicial-polytopal realization, palindromicity, and unimodality}\par
\vspace{9pt}
{\color{Muted}Research manuscript prepared for Vladimir Reshetnikov\\
20 September 2026}
\end{flushleft}
\vspace{3pt}
\hrule
\vspace{12pt}
\begin{abstract}
For a finite preorder $\tau$ on $[n]$, let $h_\tau(t)$ count the integer points of
its preorder polytope by the number of nonzero coordinates. We give a proposed
general proof that this polynomial is the $h$-polynomial of an $n$-dimensional
simplicial polytope, establishing Conjecture~5.1(c), and hence (a) and (b), of
Athanasiadis and Chapoton. The principal construction is the explicit
$n$-dimensional lattice polytope
\[
 C_\tau=\conv\bigl(\{\pm\ee_i:1\leq i\leq n\}
 \cup\{\ee_i-\ee_j:j\leq_\tau i,\ i\ne j\}\bigr).
\]
We prove $h^*(C_\tau,t)=h_\tau(t)$ by an integer-fiber decomposition of an
augmented bipartite root polytope. Every boundary pulling triangulation of
$C_\tau$, coned from the origin, is unimodular; its boundary is combinatorially
polytopal. We also prove reflexivity, an exact formula for the gauge of $C_\tau$,
and a Gorenstein translation formula for the larger root polytope. The
higher-dimensional root-polytope identity and preorder duality are already
results of Dai, Hou, Liu, Thawinrak, and Wang, and are not claimed as new here.
A separate hypertree argument reconstructs that known identity. Exact
verification code and complete small-case data accompany the manuscript.
\end{abstract}
\begin{tcolorbox}[colback=Pale,colframe=Accent,boxrule=0.5pt,arc=1mm,
  left=9pt,right=9pt,top=7pt,bottom=7pt]
\small\textbf{Scope and status.} The intended new result is the general
simplicial-polytopal realization, not the already established duality theorem.
The proof uses explicitly identified results on support-enumerators and pulling
triangulations. All steps specific to the construction are proved below; the external results
are identified explicitly. The manuscript has not been independently refereed
or checked by a proof assistant. The finite
computations test the constructions; they do not replace the proof. Flag
realizability, $\gamma$-positivity, and real-rootedness are not proved.
\end{tcolorbox}
\newpage
\tableofcontents
\vspace{8pt}
\begin{tcolorbox}[colback=Pale,colframe=Pale,arc=1mm]
\small\textbf{Reading route.} Sections~\ref{sec:problem}--\ref{sec:polytopal}
contain the main argument. Section~\ref{sec:second} gives a second geometric
check. Sections~\ref{sec:examples}--\ref{sec:limits} provide examples, exact
experiments, and limitations. Appendix~\ref{app:hypertrees} gives an alternative
proof of the known higher-dimensional identity; Appendix~\ref{app:audit}
records the proof dependencies and the main normalization checks.
\end{tcolorbox}
\newpage

\section{The problem and the result}\label{sec:problem}

A \emph{preorder} is a reflexive, transitive binary relation; antisymmetry is not
required. For a preorder $\tau$ on $E=[n]$, write $\J$ for its order ideals,
including $\varnothing$ and $E$. For a vector $x$ and a subset $A\subseteq E$,
write $x(A)=\sum_{i\in A}x_i$. The preorder polytope and its support polynomial
are
\begin{align}
 \Qt&=\{x\in\R_{\geq0}^{n}:x(I)\leq |I|\text{ for every }I\in\J\},
 \label{eq:Q}\tag{1.1}\\
 \Pt&=\Qt\cap\Z^n,\qquad
 h_\tau(t)=\sum_{x\in\Pt}t^{|\supp(x)|}
          =\sum_{k=0}^{n}h_k(\tau)t^k.\label{eq:h}
\end{align}
Here $\supp(x)=\{i:x_i\ne0\}$. These are the objects of
Athanasiadis and Chapoton~\cite{AC}. Their Conjecture~5.1(c) asks whether
$(h_0(\tau),\ldots,h_n(\tau))$ is the $h$-vector of an $n$-dimensional
simplicial polytope. Parts (a) and (b) ask, respectively, for palindromicity and
palindromic unimodality. The stronger part (d) requires a flag simplicial
polytope.

\subsection{What was already known when this problem was selected}

The arXiv version of~\cite{AC} is dated 26 May 2026. A subsequent paper by
Dai, Hou, Liu, Thawinrak, and Wang~\cite{DHLTW}, in its 27 August 2026 revision,
proves that the support-enumerator of a bipartite graph is the Ehrhart numerator
of its augmented root polytope. Its Theorem~1.3 explicitly settles
Conjecture~5.4 of~\cite{AC}, the identity $h_\tau=h_{\tau^*}$ under reversal of
the preorder. These facts must therefore be treated as prior results, not as
open problems. Problem~5.3 of~\cite{DHLTW} still discusses the coefficient-shape
conjectures for preorders. Menon's August 2026 paper~\cite{Menon} gives further
special cases of $\gamma$-positivity for arbors.

The target of the present manuscript is the \emph{general} assertion in
Conjecture~5.1(c). The advance proposed here is a dimension-reducing lattice
fiber decomposition and the resulting explicit root-polytope realization. The
literature check establishes these version-specific statements; it does not
certify absolute priority against every unpublished or unindexed argument.

\begin{theorem}[Main theorem]\label{thm:main}
Let $\tau$ be any preorder on $[n]$, where $n\geq1$, and put
\begin{equation}
 C_\tau=\conv\bigl(\{\pm\ee_i:i\in[n]\}
      \cup\{\ee_i-\ee_j:j\leq_\tau i,\ i\ne j\}\bigr)
      \subset\R^n.\label{eq:C}
\end{equation}
Then $C_\tau$ is reflexive and
\begin{equation}
 h^*(C_\tau,t)=h_\tau(t).\label{eq:main}
\end{equation}
Moreover, any pulling triangulation of $\partial C_\tau$ is combinatorially the
boundary of an $n$-dimensional simplicial polytope $S_\tau$ satisfying
\[
 h(S_\tau,t)=h_\tau(t).
\]
Consequently $h_k(\tau)=h_{n-k}(\tau)$ for all $k$, and
\[
 1=h_0(\tau)\leq h_1(\tau)\leq\cdots\leq h_{\lfloor n/2\rfloor}(\tau).
\]
In fact the full $g$-theorem conditions hold for $h(\tau)$.
\end{theorem}

A pulling triangulation is obtained by fixing an order on the vertices and
recursively triangulating faces by cones from their first vertex. The actual
root polytope $C_\tau$ need not be simplicial. The theorem asserts that its
triangulated boundary has a simplicial convex realization with the desired
$h$-vector. These two assertions should not be conflated.

The empty preorder is harmless: its support polynomial is $1$, and both root
models degenerate to a point. We treat $n\geq1$ below to avoid unnecessary
zero-dimensional boundary conventions.

\subsection{A few useful sanity checks}

Since $E$ is an ideal, every $x\in\Pt$ has $x(E)\leq n$. Hence $\Pt$ is finite,
$h_0=1$, and $h_n=1$: the only point with all coordinates positive is $\one$.
For a point supported just at $i$, the largest allowed value of its coordinate
is $|\down i|$, where $\down i=\{j:j\leq_\tau i\}$. Therefore
\begin{equation}
 h_1(\tau)=\sum_{i=1}^n|\down i|
          =|\{(j,i):j\leq_\tau i\}|.\label{eq:h1}
\end{equation}
Elements in the same equivalence class of the preorder remain separate
coordinates throughout. Replacing a preorder by its quotient poset without
retaining class sizes would change the problem.

\section{The augmented bipartite root polytope}\label{sec:bipartite}

Let $G_\tau$ have left vertices $1_L,\ldots,n_L$, right vertices
$1_R,\ldots,n_R$, and an edge $i_Lj_R$ precisely when $j\leq_\tau i$.
Adjoin $0_L$ and $0_R$, with $0_L$ adjacent to every right vertex and $0_R$
adjacent to every left vertex. Denote the resulting graph by $\widehat G_\tau$.
In particular $0_L0_R$ is an edge. Its root polytope is
\begin{equation}
 B_\tau=\conv\{\ee_i+\ff_j:i_Lj_R\in E(\widehat G_\tau)\}
       \subset\R^{n+1}\oplus\R^{n+1}.\label{eq:B}
\end{equation}
The two coordinate blocks will be denoted $(u,v)$. We use the plus-sign
convention for bipartite roots. Negating the second block converts this into
the difference convention and preserves the lattice.

The affine span of $B_\tau$ is
\begin{equation}
 H_1=\{(u,v):\textstyle\sum_{i=0}^{n}u_i
                  =\sum_{j=0}^{n}v_j=1\}.\label{eq:affine}
\end{equation}
Indeed, the edges $i_L0_R$ and $0_Lj_R$ form a spanning tree, whose $2n+1$
root vertices are affinely independent. Thus $\dim B_\tau=2n$. At dilation
$m$, the lattice is the full integer lattice in
$H_m=\{\sum u_i=\sum v_j=m\}$, not a smaller lattice generated by selected
coordinates.

\subsection{The known support-enumerator theorem}

For a bipartite graph with left neighborhoods $D_1,\ldots,D_n$ in an
$n$-element right shore, its draconian vectors are the nonnegative integer
vectors $x$ satisfying
\begin{equation}
 x(A)\leq\left|\bigcup_{i\in A}D_i\right|
       \quad(A\subseteq[n]).\label{eq:draconian}
\end{equation}
Theorem~1.2 of~\cite{DHLTW} identifies their support-enumerator with the
$h^*$-polynomial of the graph augmented by the two universal vertices just
used in~\eqref{eq:B}.

\begin{lemma}\label{lem:draconian}
For $D_i=\down i$, the draconian vectors in~\eqref{eq:draconian} are exactly
$\Pt$.
\end{lemma}
\begin{proof}
If $x\in\Pt$ and $A\subseteq E$, transitivity makes $\down A$ an ideal.
Nonnegativity gives
\[
 x(A)\leq x(\down A)\leq|\down A|
        =\left|\bigcup_{i\in A}D_i\right|.
\]
Conversely, if~\eqref{eq:draconian} holds and $I$ is an ideal, reflexivity and
downward closure give $\bigcup_{i\in I}D_i=I$. Taking $A=I$ gives
$x(I)\leq|I|$.
\end{proof}

\begin{proposition}[Known higher-dimensional identity]\label{prop:known}
For every preorder on $[n]$,
\begin{equation}
 \Ehr_{B_\tau}(t):=\sum_{m\geq0}|mB_\tau\cap\Z^{2n+2}|t^m
       =\frac{h_\tau(t)}{(1-t)^{2n+1}}.\label{eq:known}
\end{equation}
\end{proposition}
\begin{proof}
Apply the support-enumerator theorem of~\cite{DHLTW} and
Lemma~\ref{lem:draconian}.
\end{proof}

Appendix~\ref{app:hypertrees} supplies another proof of this proposition using
the K\'alm\'an--Postnikov theorem~\cite{KP}. Its elementary matching construction
also explains why a support statistic, rather than a sum-of-coordinates
statistic, occurs. Neither proof makes the already known proposition a novelty
claim of this manuscript.

\section{A complete inequality description}\label{sec:hall}

The special feature of a preorder is that all relevant neighborhood sets can
be reduced to ideals. We need the reduction over the reals, not only for
integer degree vectors.

\begin{lemma}[Weighted Hall criterion]\label{lem:flow}
Let $G$ be a bipartite graph with shores $L,R$. Nonnegative margins $u$ and $v$
with common total $m$ are the row and column sums of a nonnegative matrix
supported on the edges of $G$ if and only if
\begin{equation}
 u(A)\leq v(N(A))\qquad(A\subseteq L).\label{eq:hall}
\end{equation}
If the margins are integral, the matrix can be chosen integral.
\end{lemma}
\begin{proof}
Necessity follows by summing the entries in the rows of $A$. For sufficiency,
make a network with source-to-left capacities $u_i$, right-to-sink capacities
$v_j$, and capacity $m+1$ on each allowed left-to-right edge. A cut of capacity
at most $m$ cannot cut an allowed left-to-right edge. If its source side
contains the left set $A$, it must therefore contain $N(A)$ on the right; its
capacity is at least $m-u(A)+v(N(A))\geq m$. A flow of value $m$ follows from
the max-flow/min-cut argument. For completeness, that argument chooses a
maximum feasible flow, which exists by compactness; if its residual graph
contains a source-to-sink path, a positive augmentation contradicts maximality.
Otherwise the vertices reachable from the source give a cut with capacity equal
to the flow value. Integral capacities allow integral augmentations from the
zero flow, so an integral maximum flow exists in the integral case.
\end{proof}

\begin{proposition}[Ideal inequalities for $B_\tau$]\label{prop:inequalities}
For every real $m\geq0$, a pair $(u,v)\in H_m$ belongs to $mB_\tau$ exactly
when all its coordinates are nonnegative and
\begin{equation}
 u(I)-v(I)\leq v_0\qquad(I\in\J).\label{eq:idealB}
\end{equation}
Here $u(I),v(I)$ involve only the coordinates indexed by $[n]$.
\end{proposition}
\begin{proof}
A point of $mB_\tau$ is precisely the pair of margins of a nonnegative matrix
of total $m$ supported on $\widehat G_\tau$. For a left subset containing
$0_L$, its neighborhood is the entire right shore, so~\eqref{eq:hall} follows
from nonnegativity and the total-sum condition. For $A\subseteq[n]$,
\[
 N(A)=\{0_R\}\cup(\down A)_R.
\]
Thus the remaining Hall conditions are
$u(A)\leq v_0+v(\down A)$. If~\eqref{eq:idealB} holds, then
\[
 u(A)\leq u(\down A)\leq v_0+v(\down A),
\]
so all Hall conditions hold. Conversely, taking $A=I$ for an ideal gives
$\down I=I$ and recovers~\eqref{eq:idealB}. The case $m=0$ is immediate.
\end{proof}

The empty-ideal inequality is just $v_0\geq0$, and the full-ideal inequality
is equivalent to $u_0\geq0$. They may be redundant, but retaining them makes
later formulas uniform. We do not assert that every inequality in this
presentation defines a distinct facet.

\section{The dimension-reducing lattice fibers}\label{sec:fibers}

Define the projection
\begin{equation}
 \pi(u,v)=(u_i-v_i)_{i=1}^{n}.\label{eq:projection}
\end{equation}
On the root vertices in~\eqref{eq:B}, it has values
$0$, $\pm\ee_i$, and $\ee_i-\ee_j$ for $j\leq_\tau i$. Therefore
\begin{equation}
 \pi(B_\tau)=C_\tau.\label{eq:projectC}
\end{equation}
Its kernel on the direction space of $H_1$ is spanned by
$\ee_i+\ff_i-\ee_0-\ff_0$, $1\leq i\leq n$. The dimension loss is exactly
$n$. A linear projection alone does \emph{not} usually preserve an Ehrhart
numerator. The following explicit fibers are what make this one work.

For $y\in\R^n$, put
\begin{align}
 y_i^+&=\max(y_i,0),&y_i^-&=\max(-y_i,0),\notag\\
 p(y)&=\sum_i y_i^+,&q(y)&=\sum_i y_i^-,\label{eq:pq}\\
 a_\tau(y)&=\max_{I\in\J}y(I),&
 g_\tau(y)&=q(y)+a_\tau(y).\label{eq:gauge}
\end{align}
The empty and full ideals imply
$a_\tau(y)\geq0$ and $a_\tau(y)\geq p(y)-q(y)$.

\begin{theorem}[Exact fibers and an integral gauge]\label{thm:fibers}
For $y\in\R^n$, define
\begin{align}
 u^{\min}(y)&=(a_\tau(y)+q(y)-p(y),\ y_1^+,\ldots,y_n^+),\notag\\
 v^{\min}(y)&=(a_\tau(y),\ y_1^-,\ldots,y_n^-).\label{eq:base}
\end{align}
Every point of $mB_\tau$ projecting to $y$ has a unique expression
\begin{equation}
 (u,v)=\bigl(u^{\min}(y),v^{\min}(y)\bigr)
        +\sum_{i=0}^{n}w_i(\ee_i+\ff_i),\label{eq:fiber}
\end{equation}
where $w_i\geq0$ and
\begin{equation}
 \sum_{i=0}^{n}w_i=m-g_\tau(y).\label{eq:fibertotal}
\end{equation}
Consequently,
\begin{equation}
 y\in mC_\tau\quad\Longleftrightarrow\quad g_\tau(y)\leq m.\label{eq:gaugecriterion}
\end{equation}
For integral $y,m$, all terms in~\eqref{eq:base} are integral, and the lattice
points in the fiber correspond exactly to nonnegative integer vectors $w$
satisfying~\eqref{eq:fibertotal}.
\end{theorem}
\begin{proof}
The two vectors in~\eqref{eq:base} are nonnegative. Their sums are both
$a_\tau(y)+q(y)=g_\tau(y)$, and their difference on $[n]$ is $y$. Moreover,
\[
 u^{\min}(I)-v^{\min}(I)=y(I)\leq a_\tau(y)=v^{\min}_0.
\]
Proposition~\ref{prop:inequalities} puts this pair in $g_\tau(y)B_\tau$.
Adding a nonnegative multiple of a diagonal root leaves every difference
$u(I)-v(I)$ unchanged. For the root indexed by $0$, the right side $v_0$ also
increases, so all inequalities remain valid. This proves that every expression
in~\eqref{eq:fiber} with~\eqref{eq:fibertotal} is an admissible point.

Conversely, suppose $(u,v)\in mB_\tau$ projects to $y$. For $1\leq i\leq n$,
nonnegativity and $u_i-v_i=y_i$ imply uniquely
\[
 u_i=y_i^++w_i,\qquad v_i=y_i^-+w_i,\qquad
 w_i=\min(u_i,v_i)\geq0.
\]
The ideal inequalities give $v_0\geq\max_I y(I)=a_\tau(y)$. Set
$w_0=v_0-a_\tau(y)\geq0$. The equality of the two total sums then forces
\[
 u_0=v_0+q(y)-p(y)=u^{\min}_0(y)+w_0.
\]
All coordinates in~\eqref{eq:fiber} are now determined. Summing either block
yields~\eqref{eq:fibertotal}. Projection proves~\eqref{eq:gaugecriterion}.
If $y$ is integral, the maximum defining $a_\tau(y)$ is an integer, which
establishes the final assertion without any lattice-saturation assumption.
\end{proof}

\begin{corollary}[Equality of Ehrhart numerators]\label{cor:numerators}
For every finite preorder,
\begin{equation}
 \Ehr_{B_\tau}(t)=\frac{\Ehr_{C_\tau}(t)}{(1-t)^n},
 \qquad h^*(B_\tau,t)=h^*(C_\tau,t)=h_\tau(t).\label{eq:seriesfactor}
\end{equation}
\end{corollary}
\begin{proof}
For integral $y$, the least admissible integral dilation is $g_\tau(y)$.
Theorem~\ref{thm:fibers} gives
\[
 \Ehr_{B_\tau}(t)
 =\frac{\displaystyle\sum_{y\in\Z^n}t^{g_\tau(y)}}{(1-t)^{n+1}},
 \qquad
 \Ehr_{C_\tau}(t)
 =\frac{\displaystyle\sum_{y\in\Z^n}t^{g_\tau(y)}}{1-t}.
\]
These are legitimate formal power series: each coefficient receives only
finitely many contributions because $C_\tau$ is bounded. Dividing the two
identities proves the first equality in~\eqref{eq:seriesfactor}. The dimensions
are $2n$ and $n$, respectively, so their Ehrhart numerators coincide. Finally
apply Proposition~\ref{prop:known}.
\end{proof}

There is a useful coefficient-level form. The fiber over an integral $y$ at
dilation $m$ contains
\begin{equation}
 \binom{m-g_\tau(y)+n}{n}\quad\text{if }g_\tau(y)\leq m,
 \qquad 0\quad\text{otherwise}.\label{eq:fibercount}
\end{equation}
Thus the factor $(1-t)^{-n}$ is not a heuristic cancellation of dimensions; it
comes from a concrete, uniquely parametrized family of integer fibers.

\subsection{An order-theoretic interpretation of the gauge}

The gauge can also be written
\begin{equation}
 g_\tau(y)=\max_{I\in\J}
 \left(\sum_{i\in I}y_i^+ +\sum_{i\notin I}y_i^-\right).\label{eq:gauge2}
\end{equation}
Thus testing membership in a dilation of $C_\tau$ is a maximum-weight-ideal
problem. One need not enumerate ideals in a practical implementation: put
source capacity $y_i^+$ at $i$, sink capacity $y_i^-$ at $i$, and sufficiently
large capacities on arcs $i\to j$ for $j\leq_\tau i$. A finite minimum cut has
a downward-closed source side $I$, and its capacity is $p(y)-y(I)$. Hence
$a_\tau(y)$, and therefore $g_\tau(y)$, can be obtained by a network-flow
calculation. The accompanying small-case verifier deliberately uses exhaustive
ideal enumeration instead, keeping the check independent of a flow library.

\section{Reflexivity and simplicial-polytopal realization}\label{sec:polytopal}

We now finish the proof of Theorem~\ref{thm:main}. The required geometry is
particularly simple because all nonzero vertices of $C_\tau$ are type-$A$
roots with one coordinate deleted.

\subsection{The polar is integral}

Since $C_\tau$ contains $\conv\{\pm\ee_i\}$, it is full-dimensional and contains
$0$ in its interior. Its polar is
\begin{equation}
 C_\tau^\circ=\left\{z\in\R^n:
 -1\leq z_i\leq1,\quad z_i-z_j\leq1\text{ whenever }j\leq_\tau i\right\}.
 \label{eq:polar}
\end{equation}

\begin{lemma}[Root determinants]\label{lem:determinant}
Put $\ee_0=0\in\Z^n$. Any $n$ linearly independent vectors of the form
$\ee_i-\ee_j$, with $i,j\in\{0,\ldots,n\}$ and $i\ne j$, form a basis of
$\Z^n$.
\end{lemma}
\begin{proof}
Associate the vector $\ee_i-\ee_j$ to a directed edge from $j$ to $i$ on
$\{0,\ldots,n\}$. A cycle in the underlying undirected graph would give a
linear dependence, with signs determined by the edge directions. Consequently
$n$ independent vectors give a forest with $n$ edges on $n+1$ vertices, hence a
spanning tree. The matrix of these vectors is a directed incidence matrix with
the row indexed by $0$ deleted. Remove a leaf different from $0$; its row has
one nonzero entry, equal to $1$ or $-1$. Expanding the determinant along that
row and repeating reduces the determinant to $1$ or $-1$.
\end{proof}

\begin{proposition}\label{prop:reflexive}
The polar~\eqref{eq:polar} is a lattice polytope. In fact all its vertices have
coordinates in $\{-1,0,1\}$. Thus $C_\tau$ is reflexive.
\end{proposition}
\begin{proof}
At a vertex of the full-dimensional polar, choose $n$ linearly independent
active constraints. With $z_0=0$, each selected equality is of the form
$z_i-z_j=1$. Its coefficient row is a root as in
Lemma~\ref{lem:determinant}. The resulting square matrix has determinant
$\pm1$, so its solution is integral. The bounds $-1\leq z_i\leq1$ restrict
these integers to $-1,0,1$. Both the original polytope and its polar are
therefore integral, with the origin in the interior, which is precisely
reflexivity.
\end{proof}

This proof also supplies a finite exact way to test lattice-point membership:
one may evaluate against all ternary vectors satisfying~\eqref{eq:polar},
since that set contains every vertex of the polar. The verifier uses this
method as a check independent of the gauge formula.

\subsection{Unimodular cones over the boundary}

Recall that a lattice simplex with vertices $0,r_1,\ldots,r_n$ is unimodular
when $r_1,\ldots,r_n$ form a lattice basis. A triangulation is unimodular if all
its full-dimensional simplices are unimodular.

\begin{proposition}\label{prop:startriangulation}
Let $\Delta$ be any triangulation of $\partial C_\tau$ using only vertices of
$C_\tau$. Then the star triangulation $0*\Delta$ of $C_\tau$ is unimodular.
\end{proposition}
\begin{proof}
Every boundary simplex is contained in a supporting hyperplane not containing
$0$. The $n$ vertices of a maximal boundary simplex are therefore linearly
independent: an affine dependence is already excluded by being a simplex,
and a linear dependence with nonzero coefficient sum would place $0$ in that
supporting hyperplane. All its vertices are roots
$\ee_i-\ee_j$, allowing index $0$. Lemma~\ref{lem:determinant} shows that the
cone simplex from $0$ is unimodular. Since $0$ is interior, these cones form a
triangulation of all of $C_\tau$.
\end{proof}

\subsection{Why the triangulated boundary is polytopal}

We use the following standard pulling fact, stated as Lemma~2.1
in Athanasiadis~\cite{AthSpecial}: a pulling (reverse lexicographic)
triangulation of the boundary of a convex $n$-polytope is combinatorially the
boundary of a simplicial convex $n$-polytope. The geometric mechanism is
successive small outward perturbations of the ordered vertices; at each step
one chooses the perturbation small enough not to change the already fixed
incidences. The resulting convex boundary realizes the prescribed pulling
subdivision. This is an actual polytopality assertion, not merely shellability
or the assertion that the complex is a sphere.

Here is the Ehrhart connection, with its normalization made explicit. If $K$
is a $d$-dimensional unimodular triangulation of a lattice polytope $P$, then
\begin{equation}
 \Ehr_P(t)=\sum_{\sigma\in K}
     \left(\frac{t}{1-t}\right)^{|\sigma|}
   =\frac{h(K,t)}{(1-t)^{d+1}},\label{eq:triangEhr}
\end{equation}
where the empty face contributes $1$. Indeed, the relative interior of an
$r$-dimensional unimodular simplex contributes $\binom{m-1}{r}$ lattice points
at positive dilation $m$, with generating series
$t^{r+1}/(1-t)^{r+1}$. Relative interiors partition the triangulation, giving
the first equality. The second equality is the definition of its
$h$-polynomial. For a cone, the face generating function gains a factor
$1+t/(1-t)$, while the denominator exponent increases by one; hence
\begin{equation}
 h(0*\Delta,t)=h(\Delta,t).\label{eq:coneh}
\end{equation}

\begin{proof}[Completion of the proof of Theorem~\ref{thm:main}]
Choose a pulling triangulation $\Delta$ of $\partial C_\tau$. By the pulling
fact, $\Delta$ is the boundary of an $n$-dimensional simplicial polytope
$S_\tau$. Proposition~\ref{prop:startriangulation} and
\eqref{eq:triangEhr}--\eqref{eq:coneh} give
\[
 h(S_\tau,t)=h(\Delta,t)=h^*(C_\tau,t)=h_\tau(t),
\]
where the last equality is Corollary~\ref{cor:numerators}. Reflexivity was
proved in Proposition~\ref{prop:reflexive}. The Dehn--Sommerville relations
for simplicial polytopes give palindromicity, and the $g$-theorem gives the
stated inequalities and its stronger numerical conditions; see also the
polytopal conclusions summarized in~\cite[Theorem~3.5]{AthSpecial}.
\end{proof}

\section{A second geometric check in dimension \texorpdfstring{$2n$}{2n}}\label{sec:second}

The larger polytope $B_\tau$ supplies two independent checks on the preceding
argument: an interior-point translation and an explicit special simplex. They
are useful safeguards against a sign error or a mismatch of dimensions.

\subsection{An exact interior translation}

Let $\mathbf c=(1,\ldots,1;1,\ldots,1)\in\Z^{2n+2}$. For every $m\geq n+1$,
\begin{equation}
 (\relint(mB_\tau))\cap\Z^{2n+2}
 =\mathbf c+\bigl((m-n-1)B_\tau\cap\Z^{2n+2}\bigr).\label{eq:interior}
\end{equation}
To see this, note that every coordinate inequality is nonconstant on $B_\tau$,
and each ideal inequality has strict slack at $\ee_0+\ff_0$.
Consequently relative interior is characterized by positive coordinates and
strict versions of~\eqref{eq:idealB}. At integer points these strict
inequalities are
\[
 u(I)-v(I)\leq v_0-1.
\]
Subtracting $1$ from every coordinate reduces each block sum by $n+1$.
The terms $|I|$ cancel on the left, and the right side becomes the new
$v_0$. This is exactly the weak inequality description of
$(m-n-1)B_\tau$. Conversely, adding $\mathbf c$ makes every defining
inequality strict. There are no interior integer points for
$1\leq m<n+1$, since $n+1$ positive coordinates in either block already
require total at least $n+1$.

In particular, $B_\tau$ has codegree $n+1$, with unique interior integer point
$\mathbf c$ in its first nonempty interior dilation. It is Gorenstein of index
$n+1$. Ehrhart reciprocity translates~\eqref{eq:interior} into
\[
 t^{2n+1}h^*(B_\tau,t^{-1})=t^{n+1}h^*(B_\tau,t),
\]
which again gives $h_\tau(t)=t^nh_\tau(t^{-1})$.
This symmetry proof does not use the $g$-theorem or polytopality of a sphere.

\subsection{The diagonal special simplex}

A simplex with $r$ vertices inside a polytope is \emph{special} if each facet
contains exactly $r-1$ of those vertices. Consider
\begin{equation}
 \Sigma_\tau=\conv\{d_i:0\leq i\leq n\},
 \qquad d_i=\ee_i+\ff_i.\label{eq:diagonal}
\end{equation}
Reflexivity of the preorder makes these root vertices available, and they are
affinely independent. A coordinate hyperplane $u_i=0$ or $v_i=0$ contains
exactly the diagonal vertices other than $d_i$. On an ideal hyperplane
$u(I)-v(I)=v_0$, every $d_i$ with $i>0$ lies on the hyperplane and $d_0$ does
not. Every facet is furnished by at least one nonredundant inequality of
Proposition~\ref{prop:inequalities}; thus $\Sigma_\tau$ is special.

Every triangulation of a bipartite root polytope using its vertices is
unimodular. Indeed, affine dependence of edge roots is equivalent to a cycle
in the supporting bipartite graph; a full-dimensional simplex therefore
corresponds to a spanning tree. Leaf elimination in its incidence matrix shows
that its edge vectors form a basis of the integer lattice
$\{(u,v):\sum u_i=\sum v_j\}$. Intersecting with the height-one hyperplane
shows affine unimodularity. In particular every pulling ordering is compressed
in the terminology of~\cite{AthSpecial}.

Athanasiadis's special-simplex theorem~\cite[Theorem~3.5]{AthSpecial}, applied
with dimension $2n$ and $n+1$ special vertices, therefore realizes
$h^*(B_\tau,t)$ as the $h$-polynomial of a simplicial polytope of dimension
$2n-(n+1)+1=n$. This recovers the main theorem by another route. The broader
Gorenstein/unimodular-triangulation framework of Bruns and R\"omer~\cite{BR}
provides a further compatible interpretation, but is not needed for our main
proof.

\section{Consequences and quantitative formulas}\label{sec:consequences}

\subsection{Duality, with its prior attribution}

Swapping the coordinate blocks identifies $B_\tau$ with $B_{\tau^*}$.
This is the already established duality mechanism of~\cite{DHLTW}. In the
smaller model it becomes the especially transparent equality
\begin{equation}
 C_{\tau^*}=-C_\tau.\label{eq:dualC}
\end{equation}
Thus $h_{\tau^*}=h_\tau$. The gauge gives a third check. Ideals of $\tau^*$
are complements of ideals of $\tau$, so
\[
 a_{\tau^*}(-y)=a_\tau(y)-y(E),\qquad
 g_{\tau^*}(-y)=g_\tau(y).
\]
This rederivation is included to connect the new lower-dimensional model with
the known theorem, not to reclassify Conjecture~5.4 as open.

\subsection{Ehrhart polynomials and volumes}

The ordinary Ehrhart polynomial of $C_\tau$ is now explicit in terms of the
preorder support counts:
\begin{equation}
 |mC_\tau\cap\Z^n|=\sum_{k=0}^{n}h_k(\tau)\binom{m+n-k}{n},
 \qquad m\in\N.\label{eq:ehrC}
\end{equation}
Similarly,
\begin{equation}
 |mB_\tau\cap\Z^{2n+2}|=
 \sum_{k=0}^{n}h_k(\tau)\binom{m+2n-k}{2n}.\label{eq:ehrB}
\end{equation}
A binomial term with nonnegative upper argument smaller than its lower
argument is zero. These formulas also extend as polynomial identities in $m$.
Normalized volume, meaning dimension factorial times volume in the relevant
lattice, satisfies
\begin{equation}
 \Vol(C_\tau)=\Vol(B_\tau)=h_\tau(1)=|\Pt|.\label{eq:vol}
\end{equation}
The two volumes in this equality live in different dimensions and are both
\emph{normalized} volumes.

Remark~5.6 of~\cite{AC} considers $h_\tau(t)/(1-t)^{n+1}$. Equation~\eqref{eq:ehrC}
identifies its coefficients as the Ehrhart polynomial of $C_\tau$. This gives
that series a concrete geometric interpretation. It does not prove that its
polynomial has all complex roots on a particular vertical line; no such
root-location assertion follows just from reflexivity.

\subsection{Vertices, elementary bounds, and the $g$-vector}

Every vector displayed in~\eqref{eq:C} is a vertex. To expose $\ee_i$, choose
a linear functional taking value $1$ at coordinate $i$ and a fixed
$\varepsilon\in(0,1)$ at all other coordinates. The negative case is analogous.
To expose $\ee_i-\ee_j$, give coordinates $i,j$ the values $1,-1$ and all
others the value $0$; this root is the unique displayed vector with value $2$.
Thus
\begin{equation}
 f_0(C_\tau)=n+|\{(j,i):j\leq_\tau i\}|=n+h_1(\tau).\label{eq:vertices}
\end{equation}
Pulling adds no vertices, so the same number occurs in $S_\tau$.

The inclusions
\[
 [0,1]^n\subseteq\Qt\subseteq\{x\geq0:x(E)\leq n\}
\]
give the useful coefficientwise bounds
\begin{equation}
 \binom nk\leq h_k(\tau)\leq\binom nk^2
 \qquad(0\leq k\leq n).\label{eq:bounds}
\end{equation}
For the upper bound, choose the support in $\binom nk$ ways; the positive
integer values on it with total at most $n$ can be chosen in $\binom nk$ ways
by stars and bars. The lower bound uses the corresponding $0$--$1$ points.

For completeness, the full numerical strengthening supplied by polytopality
can be stated as follows. Put $g_0=1$ and
$g_i=h_i-h_{i-1}$ for $1\leq i\leq\lfloor n/2\rfloor$.
The $g$-theorem says this sequence is an $M$-sequence. In particular
$g_i\geq0$ and $g_{i+1}\leq g_i^{\langle i\rangle}$ whenever both sides
are relevant. Here, if
\[
 a=\binom{b_i}{i}+\binom{b_{i-1}}{i-1}+\cdots+\binom{b_j}{j},
 \qquad b_i>b_{i-1}>\cdots>b_j\geq j\geq1,
\]
is the Macaulay expansion, then
$a^{\langle i\rangle}=\binom{b_i+1}{i+1}+\cdots+\binom{b_j+1}{j+1}$,
with $0^{\langle i\rangle}=0$. These inequalities are stronger than
unimodality and follow here from an actual simplicial-polytope realization.

\section{Examples and the role of transitivity}\label{sec:examples}

\subsection{Antichains and a single equivalence class}

For an antichain, $\Qt=[0,1]^n$ and
\[
 h_\tau(t)=(1+t)^n,\qquad
 C_\tau=\conv\{\pm\ee_1,\ldots,\pm\ee_n\}.
\]
The gauge is $\sum_i|y_i|$. The smaller model is already the simplicial
cross polytope, so no boundary refinement is needed.

For the indiscrete preorder, in which every two elements are equivalent,
$\Qt=\{x\geq0:x(E)\leq n\}$. The same stars-and-bars argument used above
shows
\[
 h_\tau(t)=\sum_{k=0}^n\binom nk^2t^k.
\]
Here $B_\tau=\Delta_n\times\Delta_n$, and the lower-dimensional polytope
contains all roots $\pm\ee_i$ and $\ee_i-\ee_j$. Its gauge is
$\max(p(y),q(y))$, since its only ideals are $\varnothing,E$.
These two extreme examples check both the coefficient bounds and the change
of Ehrhart denominator.

\subsection{A non-self-dual four-element example}

Let $\tau$ have one minimum below three incomparable maximal elements. Its
polynomial, obtained directly from~\eqref{eq:Q}, is
\[
 h_\tau(t)=1+7t+12t^2+7t^3+t^4.
\]
For example, for supports of size two not containing the minimum, the three
pairs of maximal elements each allow the values $(1,1),(1,2),(2,1)$, giving
nine points; the three supports containing the minimum each allow only
$(1,1)$, giving three more. The dual order is not isomorphic to the original,
but its directed root polytope is the negative of the original one and its
support polynomial agrees. This is a useful test that duality is not merely
invariance under relabeling.

\subsection{A preorder with a nontrivial equivalence class}

Use labels $a,b,c,d,e$, with $a\sim c$, $a<e$, $c<e$, $b<e$, and $b<d$.
No other comparisons hold beyond those forced by reflexivity and transitivity.
The principal ideals are
\[
 \{a,c\},\quad\{b\},\quad\{a,c\},\quad\{b,d\},\quad\{a,b,c,e\}.
\]
Direct enumeration gives
\begin{equation}
 h_\tau(t)=1+11t+34t^2+34t^3+11t^4+t^5.\label{eq:running}
\end{equation}
This agrees with the running example of~\cite{AC}. There are $92$ points in
$\Pt$, $22$ vertices of $B_\tau$, and $16$ vertices of $C_\tau$.
The number $92$ is also the normalized volume of either root model.

For $x=(0,0,1,2,2)$, the matching construction in
Appendix~\ref{app:hypertrees} can assign the five demands to the five right
vertices as follows: the demand at $c$ uses $c_R$, the two demands at $d$ use
$b_R,d_R$, and the two demands at $e$ use $a_R,e_R$. Attach every left vertex
to $0_R$, then attach each right vertex to its assigned left vertex. This
gives the spanning-tree certificate included in the JSON data.

\subsection{Selected coefficient vectors}

\begin{center}
\begin{tabularx}{\textwidth}{@{}l r Y r@{}}
\toprule
Preorder & $n$ & $(h_0,\ldots,h_n)$ & $|P_\tau|$\\
\midrule
Antichain & 4 & $(1,4,6,4,1)$ & 16\\
Chain & 4 & $(1,10,20,10,1)$ & 42\\
Single equivalence class & 4 & $(1,16,36,16,1)$ & 70\\
One minimum, three maxima & 4 & $(1,7,12,7,1)$ & 28\\
Diamond & 4 & $(1,9,18,9,1)$ & 38\\
Example~\eqref{eq:running} & 5 & $(1,11,34,34,11,1)$ & 92\\
Six-element crown & 6 & $(1,12,48,75,48,12,1)$ & 197\\
\bottomrule
\end{tabularx}
\end{center}
The crown has minimal elements $a,b,c$ and maximal elements $A,B,C$, with
$a,b<A$, $b,c<B$, and $c,a<C$. Every row of the table is recomputed by the
supplied verifier; no floating-point root finder is involved.

\subsection{Why arbitrary reflexive relations are not enough}

Transitivity cannot simply be dropped from the fiber theorem. Consider the
neighborhoods
\[
 D_1=\{1\},\qquad D_2=\{1,2\},\qquad D_3=\{2,3\}.
\]
They define a reflexive but nontransitive relation. Its draconian inequalities
reduce to
\[
 x_i\geq0,\quad x_1\leq1,\quad x_1+x_2\leq2,\quad
 x_3\leq2,\quad x_1+x_2+x_3\leq3.
\]
The support counts are $(1,5,6,1)$: one-coordinate supports contribute
$1+2+2=5$, two-coordinate supports contribute $1+2+3=6$, and the full support
contributes just $(1,1,1)$. This is the nonpalindromic example also recorded
in~\cite[Remark~5.5]{AC}.

The general support-enumerator/root-polytope theorem remains valid for this
graph. What fails is the reduction of all weighted Hall inequalities to
matching ideal differences $u(I)-v(I)$. In particular, subtracting a common
coordinate from the two margins can damage an inequality involving different
row and column subsets. Thus the argument does not inadvertently imply
palindromicity for all graph support-enumerators.

\section{Exact computational verification}\label{sec:computations}

The program \texttt{code/verify.py} uses only Python's standard library and
integer arithmetic. It stores a preorder by bitmasks
\[
 D_i=\{j:j\leq_\tau i\}.
\]
It enumerates every reflexive relation on $[n]$ and retains precisely those
for which $j\in D_i$ implies $D_j\subseteq D_i$. This is a direct test of
transitivity. All weak compositions with total at most $n$ are then checked
against the ideal inequalities to obtain $h_\tau$.

\begin{center}
\begin{tabular}{@{}r r l@{}}
\toprule
Ground-set size & Labeled preorders checked & Result\\
\midrule
0 & 1 & Pass\\
1 & 1 & Pass\\
2 & 4 & Pass\\
3 & 29 & Pass\\
4 & 355 & Pass\\
5 & 6,942 & Pass\\
\midrule
Total & 7,332 & Pass\\
\bottomrule
\end{tabular}
\end{center}
For each of these preorders, the code checks duality, palindromicity,
unimodality, the endpoint coefficients, the relation-count formula for $h_1$,
and the special-simplex slack identities. These are finite consistency checks,
not the reason the theorem holds for arbitrary $n$.

Stronger, independent comparisons are performed for all $35$ preorders on
at most three elements. The program enumerates spanning trees directly,
deduplicates their left degree-minus-one vectors, and compares them with
$\{(n-|x|,x):x\in P_\tau\}$. It computes internal inactivity by testing all
allowed transfers to smaller coordinates, rather than replacing inactivity
by support size in the code. Across these instances it enumerates $18{,}009$
spanning trees and checks $399$ explicit matching-tree certificates.

For the same cases it counts sums of edge-root vertices directly at every
dilation $0\leq m\leq2n$. Those counts are compared both with Hall-margin
enumeration and with~\eqref{eq:ehrB}. Finite differences recover all $2n+1$
coefficients of the root-polytope Ehrhart numerator, including the predicted
zero coefficients above degree $n$. Independently, integer points of $C_\tau$
are counted using the ternary polar description, compared with the gauge
formula, and checked through $m=n$ against~\eqref{eq:ehrC}. Full-margin
fiber identities and the interior translation are also checked exactly.

Five selected four-element preorders undergo full root Ehrhart checks through
$m=8$, quotient checks through $m=4$, and direct edge-sum checks through
$m=3$. Further examples include the weighted five-element preorder and the
six-element crown. The JSON file records all individual results, and the CSV
file contains every preorder and support vector through size five.

The original conjecture paper reports tests through larger sizes for some
coefficient properties~\cite{AC}. The present computations are not offered as
an improved size record. Their purpose is to compare independent realizations
of the newly proposed construction and detect normalization, orientation, or
lattice-index errors.

\section{What remains unresolved and how to audit the claim}\label{sec:limits}

The proof gives an ordinary simplicial polytope, not a flag one. A flag complex
is determined by its edges: every collection of pairwise adjacent vertices
must form a face. Neither a pulling triangulation nor a general simplicial
polytope automatically has that property. Therefore Conjecture~5.1(d) is not
settled by the argument.

Likewise, palindromic unimodality and the $g$-theorem do not establish
$\gamma$-positivity or real-rootedness. The manuscript proves neither the
nonnegativity of all coefficients in
\[
 h_\tau(t)=\sum_{j=0}^{\lfloor n/2\rfloor}
       \gamma_jt^j(1+t)^{n-2j}
\]
nor the assertion that every zero of $h_\tau$ is real. Accordingly it does
not claim to settle Conjectures~5.2 or~5.3 of~\cite{AC}.

For independent mathematical review, the shortest audit is to verify four
links: the identification of the draconian vectors with $P_\tau$;
Proposition~\ref{prop:inequalities}; the unique integer fibers in
Theorem~\ref{thm:fibers}; and the determinant argument for cones over a
polytopal pulling boundary. Once these are accepted, the cited
support-enumerator theorem and the pulling lemma give the asserted
simplicial-polytopal realization. Section~\ref{sec:second} is a separate
geometric cross-check, not a circular use of the desired symmetry.

The result should be understood as a proposed research proof with reproducible
support, not as a claim of community acceptance. The distinction between the
known August 2026 result and the additional construction developed here is
part of the mathematical claim, not a bibliographic afterthought.

\appendix
\section{A hypertree proof of the known root identity}\label{app:hypertrees}

This appendix reconstructs Proposition~\ref{prop:known} without using the
support-enumerator theorem of~\cite{DHLTW}. It is included for conceptual
clarity and as an independent route to the same normalization. Its one
substantial imported result is the K\'alm\'an--Postnikov theorem~\cite{KP}.

\subsection{Hypertrees and internal inactivity}

For a connected bipartite graph with left shore $L$, a left hypertree is a
vector $a\in\N^L$ realized by a spanning tree $T$ through
$a_i=\deg_T(i_L)-1$. Distinct trees realizing the same vector contribute only
one hypertree. Order the left shore. A coordinate $i$ is internally inactive
when there is a smaller coordinate $j$ such that replacing $a$ by
$a-\ee_i+\ee_j$ gives another hypertree. The interior polynomial is the sum
of $t^{\text{number of internally inactive coordinates}}$ over hypertrees.

The K\'alm\'an--Postnikov theorem identifies this interior polynomial with the
Ehrhart numerator of the bipartite root polytope. Their paper uses a reversed
convention for the triangulation $h$-polynomial in its Theorem~1.1; the
Ehrhart-basis formulation immediately following that theorem, equation (1.2),
identifies the coefficients with the usual $h^*$-numerator used here. Thus
there is no additional reversal of $t$ in the statement we apply.

\begin{lemma}\label{lem:hypertrees}
The left hypertrees of $\widehat G_\tau$ are exactly
\begin{equation}
 a=(n-x(E),x_1,\ldots,x_n),\qquad x\in P_\tau.\label{eq:hypertree}
\end{equation}
\end{lemma}
\begin{proof}
A spanning tree of $\widehat G_\tau$ has $2n+1$ edges, so the sum of its left
degree-minus-one vector is $(2n+1)-(n+1)=n$. For a nonempty
$A\subseteq[n]$, the tree edges incident with $A_L$ form a forest on
$A_L\cup N(A)$. Their number is $a(A)+|A|$, and a forest on this nonempty
vertex set has at most $|A|+|N(A)|-1$ edges. Hence
\[
 a(A)\leq |N(A)|-1=|\down A|.
\]
Lemma~\ref{lem:draconian} now shows that $x=(a_1,\ldots,a_n)$ belongs to
$P_\tau$, and the sum condition determines $a_0$.

Conversely, let $x\in P_\tau$. Replace the left element $i$ by $x_i$ distinct
clones, each adjacent to the right set $\down i\subseteq[n]$. A set of clones
whose labels lie in $A$ has at most $x(A)\leq|\down A|$ members. Thus Hall's
matching condition holds, and all clones can be matched to distinct right
vertices. This is also a constructive statement: if an augmenting-path
matching leaves a clone unmatched, the alternating paths reachable from it
produce a set of clones with fewer neighbors, contradicting Hall's condition.

Build a tree by first taking all edges $i_L0_R$ for $0\leq i\leq n$. Each
right vertex $j_R$, $j>0$, is then attached to the left label of its matched
clone, or to $0_L$ if it is unmatched. The first edges form a star; every
subsequent edge attaches a new leaf. Hence the result is a spanning tree.
Its left degrees are $x_i+1$ for $i>0$ and $n-x(E)+1$ for $i=0$, exactly
as required.
\end{proof}

\begin{lemma}\label{lem:activity}
Order the left vertices with $0$ first. For the hypertree in
\eqref{eq:hypertree}, internal inactivity is $|\supp(x)|$.
\end{lemma}
\begin{proof}
The first coordinate is never internally inactive. A zero coordinate cannot
transfer a unit, so it is also not internally inactive. If $x_i>0$, then
$x-\ee_i\in P_\tau$ because all the defining inequalities are upper bounds
and the coordinates remain nonnegative. By Lemma~\ref{lem:hypertrees}, the
vector obtained by transferring one unit from coordinate $i$ to coordinate
$0$ is another hypertree. Thus every positive nonzero-index coordinate is
internally inactive, and there are no others.
\end{proof}

Summing over~\eqref{eq:hypertree} gives an interior polynomial equal to
$h_\tau(t)$. The K\'alm\'an--Postnikov theorem identifies it with
$h^*(B_\tau,t)$, proving Proposition~\ref{prop:known} by this second route.
Notice that the argument counts hypertree \emph{vectors}, not spanning trees;
counting all realizing trees would introduce incorrect multiplicities.

\section{Proof dependencies and normalization audit}\label{app:audit}

\begin{center}
\begin{tabularx}{\textwidth}{@{}p{0.29\textwidth}Y@{}}
\toprule
Claim or input & Where it comes from\\
\midrule
$h^*(B_\tau)=h_\tau$ & Known: \cite[Theorem~1.2]{DHLTW}, specialized using
Lemma~\ref{lem:draconian}. Alternatively Appendix~\ref{app:hypertrees} plus
\cite{KP}.\\[3pt]
Ideal Hall presentation & Proposition~\ref{prop:inequalities}; transitivity is
used explicitly.\\[3pt]
Integer fiber factorization & Theorem~\ref{thm:fibers}; all coordinates and
inverse parameters are displayed.\\[3pt]
$h^*(C_\tau)=h_\tau$ & Corollary~\ref{cor:numerators}; a consequence of the
fiber factorization and the known root identity.\\[3pt]
Unimodularity & Leaf elimination in Lemma~\ref{lem:determinant}, followed by
Proposition~\ref{prop:startriangulation}.\\[3pt]
Polytopal pulling boundary & The established pulling lemma
\cite[Lemma~2.1]{AthSpecial}.\\[3pt]
Palindromicity, unimodality & Dehn--Sommerville and the $g$-theorem applied to
the constructed simplicial polytope; an independent reciprocity proof of
palindromicity is in Section~\ref{sec:second}.\\[3pt]
Alternative polytopal route & The explicit special simplex and
\cite[Theorem~3.5]{AthSpecial}.\\
\bottomrule
\end{tabularx}
\end{center}

There are several easy places to make an otherwise plausible argument wrong.
The support polynomial of $Q_\tau$ is not asserted to be $h^*(Q_\tau)$.
It is the numerator of the different polytopes $B_\tau$ and $C_\tau$.
The larger denominator has exponent $2n+1$ and the smaller one exponent $n+1$.
Both auxiliary vertices, and the edge between them, are required. The edge
orientation convention is $j\leq_\tau i\iff i_Lj_R$ is present. The
integer-fiber parameters consist of $n+1$ nonnegative entries; the exponent
$n$ in~\eqref{eq:seriesfactor} results only after dividing by the additional
slack factor in the Ehrhart series of $C_\tau$.

No argument substitutes a rational equivalence for an integral one. The
fiber base is integral because $a_\tau(y)$ is an integer for integral $y$.
The determinant statements refer to the full relevant lattices. Finally,
regularity or shellability alone is not used to claim that a sphere is
polytopal: the specific pulling-boundary theorem supplies that step.

\section{Reproducing the accompanying computations}\label{app:reproduce}

The archive contains \texttt{article.tex}, \texttt{article.pdf}, a build script,
\texttt{code/verify.py}, the full CSV enumeration, a JSON result report, and
status notes. From the extracted directory, run
\begin{lstlisting}
python3 code/verify.py --output data
pdflatex -interaction=nonstopmode -halt-on-error article.tex
pdflatex -interaction=nonstopmode -halt-on-error article.tex
\end{lstlisting}
The verifier requires Python 3.10 or later and no third-party Python packages.
Do not use Python's \texttt{-O} option, since assertions are part of the checks.
The LaTeX document uses standard TeX Live packages, including
\texttt{newpxtext}, \texttt{newpxmath}, \texttt{tcolorbox}, and \texttt{hyperref}.

The CSV uses zero-based bit positions: the row value $5$ represents the set
$\{0,2\}$. Row $i$ records the principal ideal of element $i$, not its filter.
A polynomial is stored as coefficients from degree zero upward. The JSON
report also records explicit edge sets for matching certificates and complete
Ehrhart counts for the independently tested instances.

The calculations in the supplied report completed without failed assertions.
Their runtime is recorded separately as an observation from this environment,
not as a performance guarantee. No claim of formal verification is made.

\begin{thebibliography}{99}

\bibitem{AC}
C.~A. Athanasiadis and F. Chapoton,
\emph{Polytopes and posets associated to preorders},
\arxiv{2605.26916v1}, 26 May 2026.
Relevant locations: Section~5, Conjectures~5.1--5.4 and Remarks~5.5--5.6.

\bibitem{DHLTW}
Z. Dai, Q. Hou, Z. Liu, W. Thawinrak, and H. Wang,
\emph{Counting Lattice Points in Minkowski Sums of Cross Polytopes},
\arxiv{2608.16037v2}, 27 August 2026.
Relevant locations: Theorems~1.2--1.3, Sections~4.2--4.3, and Problem~5.3.

\bibitem{KP}
T. K\'alm\'an and A. Postnikov,
\emph{Root polytopes, Tutte polynomials, and a duality theorem for bipartite graphs},
Proc. London Math. Soc. \textbf{114} (2017), 561--588;
\arxiv{1602.04449}.
Relevant locations in the arXiv text: Theorem~1.1, equation~(1.2), and
Definitions~2.2--2.4.

\bibitem{AthSpecial}
C.~A. Athanasiadis,
\emph{Ehrhart polynomials, simplicial polytopes, magic squares and a conjecture of Stanley},
\arxiv{math/0312031}.
Relevant locations: Lemmas~2.1--2.2 and~3.4, and Theorem~3.5.

\bibitem{BR}
W. Bruns and T. R\"omer,
\emph{$h$-vectors of Gorenstein polytopes},
J. Combin. Theory Ser. A \textbf{114} (2007), 65--76;
\arxiv{math/0508392v2}.

\bibitem{Menon}
K. Menon,
\emph{Gamma-positivity for octopuses: a bijective proof},
\arxiv{2608.13247v1}, 13 August 2026.

\end{thebibliography}
\end{document}
